% Three views of the compliance-centre direction rule, for Section 2.7.2. The
% section builds the rule on the shortest rotation that takes the inward
% surface normal to the tool-face normal, so the figure separates the three
% quantities that construction produces:
%
%   (a) how large the mismatch is   phi, the angle between n_d = -n_s and
%                                   n_Tool, seen in an elevation whose plane is
%                                   perpendicular to the rotation axis
%   (b) where its axis lies         u_a in the surface tangent plane, located by
%                                   the direction angle alpha from +t_1
%   (c) how the displacement        r_{c,t} = ||r_{c,t}||(n_s x u_a), drawn
%       follows                     between the TCP and p_c
%
% u_a is an ordinary unit direction vector, not a rotation, so it carries one
% arrowhead and no angular components are drawn on the tangent axes. Earlier
% versions resolved a rotation vector onto t_1 and t_2, which left the reader
% unable to tell an angle from a length; that construction is gone.
%
% Because n_d and n_Tool both lie in the plane of panel (a), their cross product
% is perpendicular to it, so u_a points out of the page there. The dotted circle
% marks that direction and ties the elevation to the two plan views.
%
% Uses only tikz + arrows.meta, which config/packages.tex loads.
%
% Colour follows compliance_lever_moment.tex and case_c_direction_rule.tex: the
% surface, the tool and the frame directions are black, red carries the angular
% mismatch and the axis it defines, blue the compliance-centre displacement.
\begin{tikzpicture}[
    scale=1.1,
    font=\small,
    ax/.style={-{Latex[length=1.8mm]}, black},
    nrm/.style={-{Latex[length=2.0mm]}, thick, black},
    dir/.style={-{Latex[length=2.2mm]}, thick, red},
    lev/.style={-{Latex[length=2.2mm]}, thick, blue},
    arc/.style={-{Latex[length=1.6mm]}, red},
    surf/.style={draw=black, fill=black!8},
    tool/.style={draw=black, fill=black!4},
    lbl/.style={font=\footnotesize, inner sep=1.8pt},
    pt/.style={circle, fill=black, inner sep=0pt, minimum size=3.0pt},
    cpt/.style={circle, fill=blue, inner sep=0pt, minimum size=3.0pt},
  ]

% ============ (a) the angular mismatch between the two normals ==============
\begin{scope}[shift={(-5.05,0)}]
  \fill[surf] (-1.90,-0.22) rectangle (1.30,0);
  \draw[black] (-1.90,0) -- (1.30,0);

  \begin{scope}[shift={(-1.15,1.10)}, rotate=28]
    \draw[tool] (0,0) rectangle (1.35,0.16);
  \end{scope}

  % Both normals leave one point of the tool face, so the arc between them is
  % the angle of the rotation that takes the first to the second. The sector is
  % narrow, so the angle name sits well inside the two arrowheads, and each
  % direction name clears the face above it and the surface below it.
  \draw[nrm] (-0.709,1.335) -- (-0.709,0.155);
  \node[lbl, anchor=east] at (-0.80,0.75) {$n_d=-n_s$};
  \draw[nrm] (-0.709,1.335) -- (-0.155,0.293);
  \node[lbl, anchor=west] at (-0.19,0.94) {$n_{\mathrm{Tool}}$};

  \draw[arc] ([shift=(270:0.55)]-0.709,1.335) arc (270:298:0.55);
  \node[lbl, red] at ([shift=(284:0.80)]-0.709,1.335) {$\phi$};

  % The plane of the drawing is perpendicular to the rotation axis.
  \draw[red, fill=white] (0.72,0.62) circle (0.13);
  \fill[red] (0.72,0.62) circle (0.04);
  \node[lbl, red, anchor=west] at (0.89,0.62) {$u_a$};
\end{scope}
\node[lbl] at (-5.35,-1.55) {(a) Angular mismatch};

% ============ (b) the axis, located in the surface tangent plane ============
\begin{scope}[shift={(-2.45,0)}]
  \draw[ax] (0,0) -- (1.80,0);
  \node[lbl, anchor=west]  at (1.94,0) {$t_1$};
  \draw[ax] (0,0) -- (0,1.60);
  \node[lbl, anchor=south] at (0,1.72) {$t_2$};

  \draw[dir] (0,0) -- (35:1.10);
  \node[lbl, red, anchor=south west] at (35:1.16) {$u_a$};

  \draw[arc] (0.48,0) arc (0:35:0.48);
  \node[lbl, red] at (17.5:0.72) {$\alpha$};
\end{scope}
\node[lbl] at (-1.58,-1.55) {(b) Axis in the tangent plane};

% ============ (c) the displacement selected from that axis ==================
\begin{scope}[shift={(2.35,0)}]
  \draw[ax] (0,0) -- (1.45,0);
  \node[lbl, anchor=west]  at (1.58,0) {$t_1$};
  \draw[ax] (0,0) -- (0,1.45);
  \node[lbl, anchor=south] at (0,1.57) {$t_2$};

  \draw[black, fill=white] (-1.05,-0.75) circle (0.13);
  \fill[black] (-1.05,-0.75) circle (0.04);
  \node[lbl, anchor=east] at (-1.25,-0.75) {$n_s$};

  \draw[dir] (0,0) -- (35:1.15);
  \node[lbl, red, anchor=west] at (35:1.35) {$u_a$};

  \draw[lev] (0,0) -- (125:1.40);
  \node[cpt] at (125:1.40) {};
  \node[lbl, blue, anchor=east] at (-0.90,1.15) {$p_c$};
  \node[lbl, blue, anchor=east] at (-0.62,0.42) {$r_{c,t}$};
  \node[pt] at (0,0) {};
  \node[lbl, anchor=north west] at (0.08,-0.10) {TCP};

  % The square spans the quadrant between the two, so it is rotated by the
  % axis direction, the first of them.
  \begin{scope}[rotate=35]
    \draw[black, thin] (0.30,0) -- (0.30,0.30) -- (0,0.30);
  \end{scope}
\end{scope}
\node[lbl] at (2.43,-1.55) {(c) Selected displacement};

\end{tikzpicture}
